\Needspace{18\baselineskip}
\section{Concrete substitution, compute behaviour, and the full UAM}
\label{sec:full-uam}
This section discharges the transition-definition item in
Appendix~\ref{app:remaining-work}. All three transition families of the linear
UAM are now defined on actual typed states. Search has the contextual adequacy
proved in the preceding section. Substitution is defined by the two actual
graphs in Fig.~9(f), plugged into one focus-free simple context. Extrinsic
behaviour is a parameter constrained by the compute shape in paper \S6.4.
The assembled relation has successful final states as normal forms and gives
an instance of the abstract machine interface used earlier in this document.

Section~\ref{sec:uam-structure} proves determinism and refocusing of the full
UAM under the corresponding behaviour conditions; the concrete specimen laws
remain separate. In this unary linear substitution rule, the variable prefix
itself proves the required one-wayness, so no stability assumption on $G_S$
is needed for substitution refocusing. The definitions here use the fixed-depth typed
representation; its graph-valued presentation and depth compatibility are now
proved in Section~\ref{sec:paper-encoding}.

\subsection{Constructing the substitution fragments}
\termdef{variable-prefix}
\begin{definition}[Variable prefix]
Let $G:\star\to O$ be a hypernet. Construct $V;G$ by splitting its sole input:
add an internal term vertex $z$ and a unary variable edge from the original
input to $z$, and redirect all original source incidences at the input to $z$.
All other edges, labels, target incidences, and box bodies are retained.
\end{definition}
The implementation reuses the certified incidence construction for focus
insertion and replaces the new label by $V$. A variable and a focus have the
same singleton term source and target signatures, so the two new typing
obligations are immediate; all incidence invariants are inherited.
\lean{Graph/Substitution}{11}{26}{Linear.Substitution.prefixVariable}
\begin{lemma}[Freedom from focuses and holes]
If $G$ is focus-free, then $V;G$ is focus-free. If $G$ is hole-free, then
$V;G$ is hole-free.
\end{lemma}
\begin{proof}
Every edge is either retained from $G$ or is the new variable. In the first
case use the hypothesis, including its recursive assertions inside boxes.
The variable label is neither a focus nor a hole and has no children.
For hole freedom, extract an occurrence from a proposed hole and apply this
case distinction to its owning edge and occurrence fibre.
\end{proof}
\lean{Graph/Substitution}{29}{34}{Linear.Substitution.variable_free}
\lean{Graph/Substitution}{37}{46}{Linear.Substitution.variable_closed}
\termdef{substitution-fragments}
\begin{definition}[Substitution fragments]
For the same $G:\star\to O$, define
\[
  L(G)=\mathsf{rewrite};(V;G),\qquad R(G)=\mathsf{search};G.
\]
Both prefixes are the certified input focus insertion. These are precisely
the two shapes in Fig.~9(f), with $G$ playing the role of $G_S$ and $O$ the
ordered output list denoted by $\vec\ell$ there. No copy of $G$ is modified.
\end{definition}
\lean{Graph/Substitution}{49}{50}{Linear.Substitution.before}
\lean{Graph/Substitution}{53}{54}{Linear.Substitution.after}
\begin{lemma}[Adjacency and unique focus]
In $L(G)$ the variable source list equals the focus target list.
If $G$ is focus-free, each fragment has a unique shallow focus, with its
prescribed mode and no focus hidden in a box.
\end{lemma}
\begin{proof}
The second splitting replaces the variable's input source by the new focus
target. Both lists therefore consist of that one fresh vertex. The certified
focus-insertion theorem applies to $V;G$ using variable-prefix focus freedom,
and directly to $G$ on the right.
\end{proof}
\lean{Graph/Substitution}{57}{61}{Linear.Substitution.before_adjacent}
\lean{Graph/Substitution}{64}{67}{Linear.Substitution.before_focussed}
\lean{Graph/Substitution}{69}{71}{Linear.Substitution.after_focussed}
When $O$ is empty and $G$ is also hole-free, both fragments are actual states.
Their state constructors use the preceding focus certificates and the theorem
that focus insertion preserves hole freedom. On the left this theorem is
applied after variable-prefix hole freedom. On the right the search focus
starts at the sole input, so the resulting state is initial.
\lean{Graph/Substitution}{73}{76}{Linear.Substitution.beforeState}
\lean{Graph/Substitution}{79}{83}{Linear.Substitution.afterState}

\subsection{Context application and its consequences}
\begin{lemma}[Recognizing the unique focus]
If a labelled edge of a focussed graph has focus label $m$, it is the
distinguished focus and its mode is $m$.
\end{lemma}
\begin{proof}
A different edge would be focus-free by the focussed certificate, contradicting
its label. Having identified the edge, injectivity of the primitive focus
constructor identifies the modes. This argument works at every depth bound.
\end{proof}
\lean{Graph/Substitution}{86}{92}{Linear.Focussed.edge_of_label}
\lean{Graph/Substitution}{95}{102}{Linear.Focussed.mode_of_label}
\termdef{concrete-substitution}
\begin{definition}[Concrete substitution transition]
States $s,t$ satisfy $s\to_V t$ if there exist a closed-interface context
$C:\star\to\epsilon$, a selected shallow edge $e$ certified as its unique
hole, and a focus-free, hole-free graph $G:\star\to O$, such that the hole's
ordered signature is $\star\to O$ and
\[
 C[L(G)]\cong s,\qquad C[R(G)]\cong t.
\]
Here $C$ is focus-free, and the two substitutions use the same $C$, $e$, and
$G$. Equalities of the ordered boundary lists provide the explicit interface
transports. The isomorphisms preserve graph labels and incidences.
\end{definition}
This is application outside boxes at one place, as required in paper \S6.4.
The witnesses are concrete graphs and plugging operations, not an abstract
transition axiom. The endpoint state certificates enforce all global focussed
and hole-free restrictions.
\lean{Graph/Substitution}{106}{112}{Linear.State.Substitution}
\begin{lemma}[Modes and adjacent variable]
Every substitution starts in rewrite mode and ends in search mode. Its source
has a variable edge whose unique source is the focus target.
\end{lemma}
\begin{proof}
Transport each fragment's explicit focus label through plugging and its
endpoint isomorphism. Recognition of the unique focus proves the mode claim.
For adjacency, transport the left fragment's equality of the variable source
and focus target lists, first through boundary casts, then through the common
context attachment, and finally through the isomorphism. Recognition identifies
the transported focus with the source state's distinguished focus. Its target
list is the singleton containing the distinguished target vertex.
\end{proof}
\lean{Graph/Substitution}{115}{122}{Linear.State.substitution_modes}
\lean{Graph/Substitution}{126}{141}{Linear.State.substitution_variable}
\begin{lemma}[Endpoint renaming]
Substitution is preserved when either endpoint is replaced by a graph-isomorphic
state.
\end{lemma}
\begin{proof}
Keep the context, slot, bound computation, and boundary equalities; compose
the two witness isomorphisms with the supplied endpoint isomorphisms.
\end{proof}
\lean{Graph/Substitution}{144}{147}{Linear.State.substitution_iso}

\Needspace{18\baselineskip}
\subsection{Extrinsic behaviour and assembly}
\termdef{concrete-compute}
\begin{definition}[Compute shape]
A pair $s,t$ has compute shape if $s$ is in rewrite mode, $t$ is in search
mode, and an active operation edge of $s$ has source list consisting of the
focus target. The operation partition is a parameter. The target graph is
otherwise unrestricted beyond being a state.
\end{definition}
This follows the paper's decision to allow arbitrary effects: compute rules
need not be local graph replacements.
\lean{Graph/UAM}{11}{15}{Linear.State.Compute}
\begin{lemma}[Substitution excludes compute]
A substitution and a compute transition cannot have the same source, even if
their targets differ.
\end{lemma}
\begin{proof}
The substitution variable and the compute operation would have the same source
vertex. Source uniqueness forces their edges to coincide. Their labels would
then be equal, contradicting the distinct variable and operation constructors.
\end{proof}
\lean{Graph/UAM}{18}{24}{Linear.State.substitution_not_compute}
\begin{lemma}[Search excludes rewrite mode]
The source of a search transition is not in rewrite mode.
\end{lemma}
\begin{proof}
Extract its underlying coordinate-search witness. If the source mode were
rewrite, its starting coordinate would have rewrite mode, contradicting the
already proved absence of coordinate-search rules from that mode.
\end{proof}
\lean{Graph/UAM}{27}{37}{Linear.State.search_not_rewrite}
\termdef{uam-behaviour}
\begin{definition}[Admissible behaviour]
A behaviour $B$ is a relation on states together with a proof that every pair
it relates has compute shape. This is the paper's extrinsic parameter
$B_{\mathbb O}$. Its shape law is a requirement on the supplied parameter;
it is not an assertion that arbitrary relations satisfy that law.
Neither determinism nor locality is part of this definition.
\end{definition}
\lean{Graph/UAM}{41}{46}{Linear.Behaviour}
\termdef{full-uam}
\begin{definition}[Intrinsic, rewrite, and full transitions]
Write $\to_\bullet$ for the checked search relation and $\to_V$ for concrete
substitution. Define
\[
 \to_{\rm int}=(\to_\bullet)\cup(\to_V),\qquad
 \to_{\rm rw}=(\to_V)\cup B,\qquad
 \to=(\to_{\rm int})\cup B.
\]
\end{definition}
\lean{Graph/UAM}{49}{50}{Linear.UAM.Intrinsic}
\lean{Graph/UAM}{53}{54}{Linear.UAM.Step}
\lean{Graph/UAM}{57}{57}{Linear.UAM.Rewrite}
\begin{proposition}[Mode discipline and disjointness]
Every rewrite changes rewrite mode to search mode. Every full step is a search
or a rewrite. Search and rewrite cannot start at the same state. Intrinsic
and extrinsic transitions cannot start at the same state either.
\end{proposition}
\begin{proof}
For the rewrite-mode claim, use substitution's proved modes or the behaviour's
compute-shape certificate. Associativity of disjunction gives the search/rewrite
partition. Search excludes rewrite mode, proving the first disjointness claim.
For the second, split an intrinsic step into search or substitution: the former
is excluded by mode, and the latter by the distinct adjacent edge labels.
\end{proof}
\lean{Graph/UAM}{60}{64}{Linear.UAM.rewrite_modes}
\lean{Graph/UAM}{67}{68}{Linear.UAM.step_iff}
\lean{Graph/UAM}{71}{73}{Linear.UAM.search_rewrite_exclusive}
\lean{Graph/UAM}{76}{80}{Linear.UAM.intrinsic_extrinsic_exclusive}
\begin{theorem}[Final normality and the concrete machine]
For every admissible behaviour, successful final states have no full-UAM
successor. Hence the full transition relation and the original successful
finality predicate form a machine in the abstract sense of Section~2.
\end{theorem}
\begin{proof}
Search from a final state is impossible by the previously established search
boundary theorem. A rewrite would require rewrite mode, whereas finality
requires backtrack mode. These cases exhaust full transitions. The concrete
machine record contains exactly this relation, this finality predicate, and
this proved normality law; it carries no unproved determinism field.
\end{proof}
\lean{Graph/UAM}{83}{88}{Linear.UAM.final_normal}
\lean{Graph/UAM}{92}{95}{Linear.UAM.machine}
The empty relation is an admissible behaviour because it has no transitions
whose shape needs checking. With this parameter, a full step is equivalent to
an intrinsic step, by eliminating the false extrinsic disjunct.
\lean{Graph/UAM}{98}{99}{Linear.UAM.emptyBehaviour}
\lean{Graph/UAM}{102}{104}{Linear.UAM.empty_step}

\subsection{Actual executions and non-vacuity}
The substitution example prefixes a closed passive value by a variable and a
rewrite focus, obtaining an actual source state. A one-edge hole corolla is
its context witness. It is focus-free and has exactly its sole hole. The
corolla plugging unit isomorphisms identify the two plugged fragments with
that source and the existing initial value state. This proves an actual
substitution, not merely a conditional statement about one.
Here and elsewhere, \emph{corolla} is terminology introduced by this
formalization for a one-edge graph with its ordered boundary; it is not a term
used in the paper.
\lean{Graph/UAMExamples}{9}{10}{Linear.UAMExamples.source}
\lean{Graph/UAMExamples}{13}{29}{Linear.UAMExamples.substitute}
With empty behaviour and the passive partition, this substitution followed by
the already checked value-search transition gives exactly two full-machine
steps to a final state. The counted execution is built by two applications of
the execution constructor and its zero-step base case. The source cannot have
compute shape for any operation partition, by the general exclusion theorem.
\lean{Graph/UAMExamples}{33}{38}{Linear.UAMExamples.terminates}
\lean{Graph/UAMExamples}{41}{42}{Linear.UAMExamples.source_not_compute}
A separate example changes the operation partition to active and places a
rewrite focus before the value symbol. Directly unfolding the incidence lists
proves compute shape for the pair consisting of this state and the initial
value state. The singleton relation containing that pair is therefore an
admissible behaviour; equality substitution reduces its shape certificate to
the checked example. Its pair is an actual full-UAM step by the extrinsic
injection. This tests the behaviour parameter independently of substitution.
\lean{Graph/UAMExamples}{45}{47}{Linear.UAMExamples.activeSource}
\lean{Graph/UAMExamples}{51}{55}{Linear.UAMExamples.compute_shape}
\lean{Graph/UAMExamples}{58}{60}{Linear.UAMExamples.behaviour}
\lean{Graph/UAMExamples}{62}{63}{Linear.UAMExamples.compute_step}

The following section discharges the tuple and concrete-definition obligation.
Section~\ref{sec:uam-structure} now supplies substitution rootedness and
uniqueness and the full-UAM structural theorem. The concrete specimen
interpretation and final robustness instantiation remain unproved.
